% =====================================================================
% REUSABLE PROJECT REPORT TEMPLATE CREATED BY JOSH MACDONALD
%The write up here is related to script 1 from the github repo, Regression Models. 

%The purpose of this write up is to conclude my findings and summarise them in a professional manner. Giving me some valuable experience writing reports and expaling the code and findings.

%Feel free to reuse the report template for any of your own personal projects. The sturcture for this one is inspired by my mathematical writing workshops during my honours complex variables course.
% =====================================================================
\documentclass[11pt,a4paper]{article}

% --- Page layout ---------------------------------------
\usepackage[margin=25mm,headheight=15pt]{geometry}
\usepackage[T1]{fontenc}
\usepackage[utf8]{inputenc}
\usepackage{lmodern}
\usepackage{microtype}
\usepackage{parskip}
\usepackage{setspace}
\setstretch{1.08}
\usepackage{enumitem}
\setlist{itemsep=3pt,topsep=5pt}
\setlength{\emergencystretch}{2em}

% --- Mathematical notation, tables and figures ------------------------
\usepackage{amsmath,amssymb}
\usepackage{graphicx}
\usepackage{booktabs} 
\usepackage{tabularx}
\usepackage{caption}
\captionsetup{font=small,labelfont=bf,labelsep=period}
\newcommand{\R}{\mathbb{R}}
\newcommand{\Z}{\mathbb{Z}}
\newcommand{\N}{\mathbb{N}}

% --- Restrained colour palette ---------------------------------------
\usepackage{xcolor}
\definecolor{Navy}{HTML}{183247}
\definecolor{Teal}{HTML}{167D8D}
\definecolor{Muted}{HTML}{5D6974}
\definecolor{Pale}{HTML}{F0F5F7}

% --- Heading styles --------------------------------------------------
\usepackage{titlesec}
\titleformat{\section}
  {\Large\sffamily\bfseries\color{Navy}}{\thesection}{0.7em}{}
\titleformat{\subsection}
  {\large\sffamily\bfseries\color{Navy}}{\thesubsection}{0.7em}{}
\titleformat{\subsubsection}
  {\normalsize\sffamily\bfseries\color{Navy}}{\thesubsubsection}{0.7em}{}
\titlespacing*{\section}{0pt}{20pt}{9pt}
\titlespacing*{\subsection}{0pt}{13pt}{6pt}

% --- A reusable highlighted finding box -------------------------------
\usepackage{tcolorbox}
\newtcolorbox{findingbox}[1][Key finding]{
  colback=Pale,colframe=Teal,boxrule=0.5pt,arc=2pt,
  colbacktitle=Teal,coltitle=white,
  fonttitle=\sffamily\bfseries,title=#1,
  left=10pt,right=10pt,top=8pt,bottom=8pt
}

% --- Links and PDF navigation ----------------------------------------
\usepackage{hyperref}
\hypersetup{colorlinks=true,linkcolor=Navy,citecolor=Teal,urlcolor=Teal}
\urlstyle{same}

% =====================================================================
% MAKE SURE TO EDIT THESE DETAILS FOR EACH REPORT
% Use a short title in the page header to prevent it overflowing.
% =====================================================================
\newcommand{\ReportTitle}{Predicting Vehicle Fuel Efficiency Using Linear Regression} %Full title
\newcommand{\ReportSubtitle}{An investigation of the Auto MPG dataset} %Conicse description of title
\newcommand{\ShortTitle}{Auto MPG Investigation}
\newcommand{\AuthorName}{Josh MacDonald}
\newcommand{\Affiliation}{Independent Project}
\newcommand{\ReportDate}{\today}
\newcommand{\ReportVersion}{1.0}
\hypersetup{pdftitle={\ReportTitle},pdfauthor={\AuthorName}}

% --- Headers and footers ---------------------------------------------
\usepackage{fancyhdr}
\pagestyle{fancy}
\fancyhf{}
\fancyhead[L]{\small\sffamily\color{Muted}\ShortTitle}
\fancyfoot[L]{\small\sffamily\color{Muted}\AuthorName}
\fancyfoot[R]{\small\sffamily\color{Muted}\thepage}
\renewcommand{\headrulewidth}{0.3pt}
\renewcommand{\footrulewidth}{0pt}

\begin{document}

% --- Cover page -------------------------------------------------------
\begin{titlepage}
\thispagestyle{empty}
\vspace*{12mm}
{\sffamily\small\bfseries\color{Teal} PROJECT REPORT\par}
\vspace{7mm}
{\color{Teal}\rule{22mm}{1.5pt}\par}
\vspace{10mm}
{\sffamily\Huge\bfseries\color{Navy}\ReportTitle\par}
\vspace{6mm}
{\sffamily\Large\color{Muted}\ReportSubtitle\par}
\vspace{18mm}
\begin{minipage}{0.85\textwidth}
{\large  This report investigates how accurately linear regression 
can predict vehicle fuel efficiency in the Auto MPG dataset, and 
whether adding horsepower and model year improves predictions 
compared with using vehicle weight alone. 
The investigation explores the relationship between vehicle 
characteristics and fuel efficiency while assessing the usefulness 
and limitations of simple predictive models. \par} %Change this to change the cover page aims
\end{minipage}
\vfill
{\color{Navy}\rule{\textwidth}{0.5pt}\par}
\vspace{5mm}
\begin{tabularx}{\textwidth}{@{}X r@{}}
\textbf{\AuthorName} & \ReportDate \\
\Affiliation & Version \ReportVersion
\end{tabularx}
\vspace*{8mm}
\end{titlepage}

% --- Front matter -----------------------------------------------------
% For a short report, remove the contents page and its clearpage command.
\pagenumbering{roman} %write this last around 150 -250 words state the question, approach, main result and practical conclusion. Include the most important limitation.
\section*{Executive Summary}
\addcontentsline{toc}{section}{Executive Summary}

This investigation examines how accurately linear regression
can predict vehicle fuel efficiency in the Auto MPG dataset,
and whether horsepower and model year improve predictions
beyond vehicle weight alone.

After removing six observations with missing horsepower values,
392 complete observations remained. Three approaches were
compared: a baseline predicting the training mean, a weight-only
linear regression, and a multiple linear regression using weight,
horsepower and model year. Evaluation used five-fold
cross-validation within the training set and a separate test
set of 79 observations.

Weight-only regression reduced test mean absolute error from
5.881 to 3.464 miles per gallon, an improvement of 41.1\%
over the baseline. Multiple regression reduced this error
further to 2.509 miles per gallon, improving on weight-only
regression by 27.6\%. It also achieved the lowest mean
cross-validation error, at 2.695 miles per gallon.

These results indicate that horsepower and model year jointly
add predictive value beyond weight. However, systematic
residual patterns remained, suggesting that the linear models
did not capture all relevant variation. The evaluation concerns
vehicles from model years 1970--1982 and does not establish
accuracy for modern vehicles. Exploratory analysis also used
the full dataset, so any resulting influence on predictor
selection limits the independence of the test evaluation.

%Enviroment to add key takeaways or findings
\begin{findingbox}[Main takeaway]
Adding horsepower and model year to weight reduced test mean
absolute error by 27.6\%, from 3.464 to 2.509 miles per gallon.
The multiple regression model performed best among the three
approaches evaluated on this historical dataset.
\end{findingbox}

\clearpage
\tableofcontents
\clearpage
\pagenumbering{arabic}

% =====================================================================
% =====================================================================
\section{Introduction}
\label{sec:introduction} %Introduce the problem and and topic keep it concise

A vehicle's fuel efficiency describes the distance it can
travel on a given amount of fuel. In this dataset, it is
measured in miles per gallon (MPG), with higher values
indicating greater fuel efficiency.

In this investigation, I examine how vehicle weight,
horsepower and model year relate to fuel efficiency.
I compare a regression model using weight alone with one
that also includes horsepower and model year, to assess
whether the additional information improves predictions.

%====================================================================
%Write one Clear question, then answer the following questions Objective 1 - what will be investigated or built?
%Objective 2 -  what will be compared or measured?
%


\subsection{Research question and objectives}
\textbf{Research question:} How accurately can linear regression
predict fuel efficiency in the Auto MPG dataset, and does including
horsepower and model year improve predictions compared with using 
vehicle weight alone?

\begin{itemize}
  \item Prepare and explore the Auto MPG dataset, checking data 
  quality and examining relationships between vehicle 
  characteristics and fuel efficiency, measured in miles per gallon.
  
  \item Fit a linear regression model using vehicle weight alone and
  a multiple linear regression model using weight, horsepower and 
  model year. Compare both with a baseline that predicts the average
  fuel efficiency of vehicles in the training data.
  
  \item Evaluate predictive performance using five-fold cross-
  validation on the training data. Measure performance on the test 
  data using metrics such as mean absolute error, root mean squared error and the coefficient of determination ($R^2$). The residual plots are examined as well to identify patterns in prediction errors.
\end{itemize}


% Define what the investigation covers, its intended use and any relevant boundaries.
\subsection{Scope} 

This investigation covers data cleaning, exploratory analysis,
model fitting and evaluation using the Auto MPG dataset.
It compares two linear regression models with a baseline
that predicts the mean fuel efficiency of the training data.

I chose this project to put the data-analysis methods I have
studied into practice and gain experience explaining the
results in a report. The analysis is limited to the historical
vehicles in the dataset. Regularised and nonlinear models
are outside its scope, but could be explored in future work.

\section{Data or Materials}
%Describe the dataset, experimental materials or other inputs. State sources, units, sample sizes, date ranges and relevant restrictions.
\label{sec:data}
This investigation uses the Auto MPG dataset from the University of
California, Irvine Machine Learning Repository \cite{auto-mpg}.
The input dataset contains 398 observations and nine variables,
covering vehicles from model years 1970 to 1982. The outcome of
interest is city-cycle fuel efficiency, measured in miles per gallon.

The recorded vehicle characteristics are cylinder count, engine
displacement, horsepower, weight, acceleration, model year and origin.
A vehicle name is also included. The regression models in this
investigation use weight, measured in pounds; engine power, measured
in horsepower; and model year, recorded as a two-digit value.

The dataset provides historical observations for exploring relationships
and evaluating predictions. Its coverage does not support assuming
that the findings apply to modern vehicles.

%Describe cleaning, exclusions, missing values, transformations and quality checks, explain decisions that could affect the results.
\subsection{Preparation and quality checks}

The raw data were loaded using the pandas library in Python. Column
names were assigned during import, and question marks were interpreted
as missing values. Initial checks examined the dataset dimensions,
column types, missing values, duplicate rows and numerical summary
statistics.

Six observations contained missing horsepower values. The cleaning
procedure removed duplicate rows, if present, and excluded observations
with any missing values. This produced a dataset of 392 complete
observations, which was saved as a separate file for exploratory
analysis and modelling. Missing values were not estimated or replaced,
and the original measurement scales were retained.

The same cleaned dataset was used for both regression models, allowing
their performance to be compared using the same observations. This
also excluded the six incomplete observations from the weight-only
model, although horsepower was not required by that model. Removing
these observations simplified the comparison but could introduce bias
if the excluded vehicles differed systematically from those retained.

% Example of a table with automatic wrapping of the description column.
\begin{table}[htbp]
\centering
\caption{Outcome and predictors used in the regression models.}
\label{tab:variables}
\begin{tabularx}{\linewidth}{@{}l X l@{}}
\toprule
\textbf{Variable} & \textbf{Description} & \textbf{Unit} \\
\midrule
Fuel efficiency & City-cycle fuel efficiency; the outcome predicted
by both regression models. & Miles per gallon \\
Weight & Vehicle weight; a predictor in both regression models.
& Pounds \\
Horsepower & Engine power; an additional predictor in the multiple
regression model. & Horsepower \\
Model year & Vehicle model year, recorded from 70 to 82 to represent
1970 to 1982; an additional predictor in the multiple regression
model. & Year (two-digit code) \\
\bottomrule
\end{tabularx}
\end{table}


\section{Methods}
%[Explain the approach in enough detail for another person to reproduce it. Justify the important choices and state relevant assumptions
\label{sec:methods}

The analysis used the 392 complete observations obtained after
cleaning. Fuel efficiency was treated as a continuous outcome.
Three approaches were compared: a baseline predicting the mean
fuel efficiency of the training observations, a simple linear
regression using vehicle weight, and a multiple linear regression
using weight, horsepower and model year.

The data were divided into training and test sets using an
80:20 split with a random seed of 42. This produced 313 training
observations and 79 test observations. Model coefficients were
estimated using the training observations only. Five-fold
cross-validation within the training set was used to assess
predictive performance, followed by evaluation on the test set.

%Document software version, random seeds, commit used and code versions for reproducibility- add a repo link where appropriate
\subsection{Reproducibility}
\label{sec:reproducibility}

The analysis was implemented in Python using pandas for data
handling, Matplotlib and Seaborn for visualisation, and
scikit-learn for regression modelling and evaluation.
The code and data are available in the project repository:
\url{https://github.com/joshmacd/Regression-Models}.

The workflow comprises four scripts, executed in the following
order:
\begin{enumerate}
    \item \texttt{01\_clean\_data.py}: inspect and clean the data.
    \item \texttt{02\_eda.py}: explore the data and generate figures.
    \item \texttt{03\_simple\_linear\_regression.py}: fit and
    evaluate the weight-only model and mean baseline.
    \item \texttt{04\_multiple\_linear\_regression.py}: fit the
    multiple regression model and compare all three approaches.
\end{enumerate}
These scripts are located in \texttt{scripts/auto\_mpg/}
and should be run from the repository's root directory.

The training and test split used
\texttt{test\_size=0.20} and \texttt{random\_state=42}.
Cross-validation used five folds with shuffling enabled and
\texttt{random\_state=42}. Both regression models included
an intercept and used unscaled predictors. The baseline used
\texttt{DummyRegressor(strategy="mean")}, and cross-validation
used \texttt{scoring="neg\_mean\_absolute\_error"}.
The returned scores were multiplied by $-1$ to report positive
mean absolute errors.

The results reported here were generated the software
versions listed in Table~\ref{tab:software}.

%Check  on VSCODE, using by checking the python -m pip and python --versions
\begin{table}[htbp]
\centering
\caption{Software versions used to generate the reported results.}
\label{tab:software}
\begin{tabular}{@{}ll@{}}
\toprule
\textbf{Software} & \textbf{Version} \\
\midrule
Python & 3.12.1 \\
pandas & 3.0.5 \\
Matplotlib & 3.11.1 \\
Seaborn & 0.13.2 \\
scikit-learn & 1.9.0 \\
NumPy & 2.5.2 \\
SciPy & 1.18.1 \\
\bottomrule
\end{tabular}
\end{table}


%Describe the method, try to be sufficiently rigourous.
\subsection{Model or analytical approach}
%Adding a note to state that the mathematics behind the models in not totally relevant and may be skipped.
\begin{quote}
\begin{quote}
\textit{I have included the following derivations for
completeness and because I enjoy the mathematics behind
the methods. Readers interested mainly in the findings
can skip the proofs without losing the main argument.}
\end{quote}
\end{quote}
To begin, let $y_i \in \R $ denote the observed fuel efficiency of
vehicle $i$, and let $w_i$, $h_i \in \R$ and $t_i \in \N$ denote its
weight, horsepower and model year, respectively. The two regression
specifications were:

\begin{align}
    y_i &= \beta_0 + \beta_1 w_i + \varepsilon_i,
    \label{eq:simple-model}\\
    y_i &= \beta_0 + \beta_1 w_i + \beta_2 h_i
          + \beta_3 t_i + \varepsilon_i,
    \label{eq:multiple-model}
\end{align}

where $\varepsilon_i$ represents variation not captured by the
specified linear relationship. The coefficients are estimated
separately for each model.

In the multiple regression model, each slope describes the
change in predicted fuel efficiency associated with a one-unit
increase in its predictor, holding the other predictors fixed.
These coefficients describe associations and are not interpreted
as causal effects. The intercept corresponds to all predictors
being zero, which lies outside the observed vehicle characteristics
and therefore has no useful physical interpretation here.

\subsubsection{Least-squares estimation, \cite{illinois-least-squares}.}

For either regression model, let $n$ be the number of training
observations and $p$ the number of predictors. Define the response
vector and design matrix by
\[
    \mathbf{y} =
    \begin{pmatrix}
        y_1\\
        \vdots\\
        y_n
    \end{pmatrix} \in \R^n,
    \qquad
    \mathbf{X} =
    \begin{pmatrix}
        1 & x_{11} & \cdots & x_{1p}\\
        \vdots & \vdots & & \vdots\\
        1 & x_{n1} & \cdots & x_{np}
    \end{pmatrix} \in \R^{n \times (p+1)}.
\]
The first column of $\mathbf{X}$ accounts for the intercept.
Ordinary least squares estimates the coefficient vector by
minimising the residual sum of squares:
\begin{equation}
    \widehat{\boldsymbol{\beta}}
    \in \operatorname*{arg\,min}_{\mathbf{b}\in\mathbb{R}^{p+1}}
    S(\mathbf{b}),
    \qquad
    S(\mathbf{b})
    = \|\mathbf{y}-\mathbf{X}\mathbf{b}\|_2^2.
    \label{eq:least-squares}
\end{equation}

\textbf{Proposition.}
If $\mathbf{X}$ has linearly independent columns, the unique
least-squares solution is
\begin{equation}
    \widehat{\boldsymbol{\beta}}
    = (\mathbf{X}^{\mathsf T}\mathbf{X})^{-1}
      \mathbf{X}^{\mathsf T}\mathbf{y}.
    \label{eq:ols-estimator}
\end{equation}

\textit{Proof.}
Expanding the objective gives
\[
    S(\mathbf{b}) = \|\mathbf{y}-\mathbf{X}\mathbf{b}\|_2^2= (\mathbf{y}-\mathbf{X}\mathbf{b})^{\mathsf T}(\mathbf{y} -  \mathbf{X}\mathbf{b})
    = \mathbf{y}^{\mathsf T}\mathbf{y}
      - 2\mathbf{b}^{\mathsf T}\mathbf{X}^{\mathsf T}\mathbf{y}
      + \mathbf{b}^{\mathsf T}\mathbf{X}^{\mathsf T}
        \mathbf{X}\mathbf{b}.
\]
Differentiation with respect to $\mathbf{b}$ yields
\[
    \frac{d S(\mathbf{b})}{d \mathbf{b}}
    = -2\mathbf{X}^{\mathsf T}\mathbf{y}
      +2\mathbf{X}^{\mathsf T}\mathbf{X}\mathbf{b}.
\]
A stationary point therefore satisfies the following
\[
    \frac{d S(\mathbf{b})}{d \mathbf{b}}
    = -2\mathbf{X}^{\mathsf T}\mathbf{y}
      +2\mathbf{X}^{\mathsf T}\mathbf{X}\mathbf{b} = 0.
\]
Hence we have
\[
    \mathbf{X}^{\mathsf T}\mathbf{X}
    \widehat{\boldsymbol{\beta}}
    = \mathbf{X}^{\mathsf T}\mathbf{y}.
\]
For every nonzero vector $\mathbf{v}$, full column rank implies that $\mathbf{X}\mathbf{v} \neq 0$. By positive definiteness of the Euclidean norm,
\[
    \mathbf{v}^{\mathsf T}\mathbf{X}^{\mathsf T}
    \mathbf{X}\mathbf{v}
    = \|\mathbf{X}\mathbf{v}\|_2^2 > 0.
\]
Thus $\mathbf{X}^{\mathsf T}\mathbf{X}$ is positive definite and hence invertible. The Hessian of $S$ is
\[
2\mathbf{X}^{\mathsf T}\mathbf{X},
\]
which is also positive definite, so the stationary point is the unique global minimum.
\hfill $\square$

Equation~\ref{eq:ols-estimator} is a theoretical expression.
The models were fitted using the \texttt{LinearRegression}
estimator in scikit-learn, rather than explicitly calculating
the matrix inverse. Full column rank is the condition for the
uniqueness result above; it is not established by the proof
for this particular dataset.

Predictions for an observation with predictor vector
$\mathbf{x}_i$ are consequently
\[
    \widehat{y}_i
    = \widehat{\beta}_0
      + \sum_{j=1}^{p}\widehat{\beta}_j x_{ij}.
\]

\subsubsection{Mean-prediction baseline}

The baseline predicts the same value for every vehicle:
the mean fuel efficiency of the training observations,
\[
    \overline{y}_{\mathrm{train}}
    = \frac{1}{n}\sum_{i=1}^{n}y_i.
\]
This is the constant prediction that minimises squared error
on the training data.

\textit{Proof - \cite[Section 2.1, Theorem 1(a)]{iowa-mean-proof}.}
For any constant $c \in \R$, we see that $y_i - c  = y_i - \overline{y}_{\mathrm{train}} + \overline{y}_{\mathrm{train}} - c$ and hence
\[
    \sum_{i=1}^{n}(y_i-c)^2
    =
    \sum_{i=1}^{n}(y_i-\overline{y}_{\mathrm{train}})^2
    +n(c-\overline{y}_{\mathrm{train}})^2.
\]
Since
\[
    \sum_{i=1}^{n}(y_i-\overline y_{\mathrm{train}})
    = \sum_{i=1}^{n}y_i-n\overline y_{\mathrm{train}}=0,
\]
the cross-term in the expansion vanishes. The first term,
$\sum_{i=1}^{n}(y_i-\overline{y}_{\mathrm{train}})^2$,
depends only on the observed training data and is therefore
constant with respect to $c$. For every $c\in\mathbb{R}$,
\[
    (c-\overline{y}_{\mathrm{train}})^2 \geq 0.
\]
Since $n>0$, it follows that
\[
    n(c-\overline{y}_{\mathrm{train}})^2 \geq 0.
\]
Hence,
\[
    \sum_{i=1}^{n}(y_i-c)^2
    =
    \sum_{i=1}^{n}(y_i-\overline{y}_{\mathrm{train}})^2
    + n(c-\overline{y}_{\mathrm{train}})^2
    \geq
    \sum_{i=1}^{n}(y_i-\overline{y}_{\mathrm{train}})^2.
\]
This establishes a lower bound on the sum of squared errors
for any constant prediction $c$.

The bound is attained if and only if
\[
    n(c-\overline{y}_{\mathrm{train}})^2=0
    \quad\Longleftrightarrow\quad
    c=\overline{y}_{\mathrm{train}}.
\]
Thus the training mean attains the lower bound, and every other
constant prediction produces a strictly larger sum of squared
errors. Therefore, $c=\overline{y}_{\mathrm{train}}$ is the unique
minimiser.
\hfill $\square$

This baseline measures the benefit of using vehicle
characteristics rather than a single constant prediction.
Its optimality concerns training squared error; it does not
imply optimality for absolute error or test-set performance.

\subsubsection{Why evaluation requires unseen observations}

A model's ability to fit the training observations does not, by
itself, establish its ability to predict fuel efficiency for
other vehicles. Additional predictors may capture useful
relationships but they may also fit random variation specific
to the training sample, this  behaviour is known as
overfitting, \cite[Section 2.2]{islr}.

The weight-only model is contained within the multiple
regression model. We have that, setting the horsepower and model-year
coefficients to zero gives the following
\[
    \beta_0+\beta_1w_i+0h_i+0t_i
    =\beta_0+\beta_1w_i.
\]
Thus every prediction for the weight-only model is
also available to the multiple regression model, \cite{penn-state-nested-models}.

\textbf{Proposition.}
When both models are fitted by ordinary least squares to the
same training observations, the minimum training sum of squared
errors for the multiple regression model cannot exceed that
of the weight-only model.

\textit{Proof.}
Define the two training sum of squared errors as
\begin{align*}
    S_{\mathrm{weight}}(a_0,a_1)
    &=
    \sum_{i=1}^{n}
    \left(y_i-a_0-a_1w_i\right)^2,\\
    S_{\mathrm{multiple}}(b_0,b_1,b_2,b_3)
    &=
    \sum_{i=1}^{n}
    \left(y_i-b_0-b_1w_i-b_2h_i-b_3t_i\right)^2.
\end{align*}
Let $(\widehat{a}_0,\widehat{a}_1)$ be a least-squares
solution for the weight-only model. The coefficient vector
$(\widehat{a}_0,\widehat{a}_1,0,0)$ is an available choice
for the multiple regression model and satisfies
\[
    S_{\mathrm{multiple}}
    (\widehat{a}_0,\widehat{a}_1,0,0)
    =
    S_{\mathrm{weight}}(\widehat{a}_0,\widehat{a}_1).
\]
Minimising over all four coefficients cannot give a larger
objective value than this particular choice. Therefore,
\begin{align*}
    \min_{(b_0,b_1,b_2,b_3)\in\mathbb{R}^4}
    S_{\mathrm{multiple}}(b_0,b_1,b_2,b_3)
    &\leq
    S_{\mathrm{multiple}}
    (\widehat{a}_0,\widehat{a}_1,0,0)\\
    &=
    \min_{(a_0,a_1)\in\mathbb{R}^2}
    S_{\mathrm{weight}}(a_0,a_1).
\end{align*}
This establishes the claimed inequality.
\hfill $\square$

However, we should note that the inequality above need not be strict: additional predictors do not necessarily reduce the training sum of squared errors.
Furthermore, the result concerns squared error on the training
observations only. It does not guarantee an improvement in
training absolute error or in performance on unseen observations.

For this reason, predictive performance was assessed using
five-fold cross-validation, \cite[Section 5.1.3]{islr}, within the training set. The training
observations were divided into five groups. Each model was
fitted using four groups and evaluated on the remaining group,
with this process repeated until every group had been used
for validation. The same folds were used for all three
approaches for a consistent comparison.

After cross-validation, each model was fitted using all 313
training observations and evaluated on the same 79 test
observations. These test observations were excluded from
coefficient estimation and cross-validation. This separates
the assessment of predictive performance from the observations
used to fit the models.

\subsubsection{Assumptions and interpretation}

Fitting a least-squares model does not show that a linear
relationship is appropriate for the data. This section
sets out the assumptions used to interpret the models
and explains what the residual plots can tell us.

\textbf{Conditional mean specification.}
We can write a general regression model in the following way
\[
    y_i
    =\beta_0+\sum_{j=1}^{p}\beta_jx_{ij}+\varepsilon_i,
\]
and here we can interpret the linear expression as the conditional mean
of fuel efficiency which assumes that
\[
    \mathbb{E}[\varepsilon_i\mid\mathbf{x}_i]=0,
\]
where $\mathbf{x}_i$ contains the predictor values for vehicle
$i$. Equivalently,
\[
    \mathbb{E}[y_i\mid\mathbf{x}_i]
    =\beta_0+\sum_{j=1}^{p}\beta_jx_{ij}.
\]
This means that, for a given combination of predictor values,
the unexplained component averages to zero in the population
being modelled. Individual observations may still differ
substantially from the conditional mean.

This assumption is not established merely by fitting a
regression model. In particular, the fitted training residuals
sum to zero when ordinary least squares includes an intercept,
but this algebraic property does not prove that the conditional
mean has been correctly specified. For background on the regression model and its error term, see \cite[Sections 3.1--3.2]{islr}.

\textbf{Linearity and additivity.}
The models used in this investigation contain separate linear
terms for the selected predictors, with no interaction or
nonlinear terms. In the multiple regression model, a one-unit
increase in predictor $j$ changes predicted fuel efficiency
by $\widehat{\beta}_j$, holding the other predictors fixed.

Take for example, the predicted change associated with a one-pound
increase in weight is the same at every horsepower and model
year. A term such as $w_i^2$ would allow curvature in the
relationship with weight, while a term such as $w_ih_i$ would
allow the weight slope to depend on horsepower. Neither type
of term was included in the fitted models. Interactions and nonlinear extensions of linear regression are discussed in \cite[Section 3.3.2]{islr}.

\textbf{Error distribution and variability.}
Calculating the least-squares fit does not require normally
distributed errors or constant error variance. Constant
conditional error variance, which is commonly known as homoscedasticity,
implies that
\[
    \operatorname{Var}(\varepsilon_i\mid\mathbf{x}_i)
    =\sigma^2
\]
for all predictor values. The spread of
the unexplained variation would remain constant across different types of vehicles.

Additional distributional and dependence assumptions are
relevant to conventional coefficient standard errors,
confidence intervals and significance tests. For the associated inferential framework, see \cite[Sections 3.1--3.2]{islr}. These
procedures were not carried out in this investigation.
Nevertheless, changing error variability remains relevant
because it may indicate that predictions are less reliable
for some vehicles than for others.

\textbf{Residual assessment.}
A residual is the difference between an observed outcome and
its fitted prediction:
\[
    r_i=y_i-\widehat{y}_i.
\]
Residuals are observable after fitting, whereas the theoretical
errors $\varepsilon_i$ depend on the unknown underlying model.

Plots of test residuals against predicted fuel efficiency were
examined for systematic patterns. Curvature may suggest that
the linear specification misses part of the relationship,
while a changing vertical spread may indicate unequal error
variability. These plots provide diagnostic evidence but do
not prove that the model assumptions hold.

\textbf{Generalisability and causal interpretation.}
Interpreting validation performance as an indication of
performance on other vehicles requires those vehicles to be
sufficiently comparable to the observed sample. Closely related
observations appearing in both training and evaluation data
could make prediction appear easier than it would be for
substantially different vehicles.

The random split assesses prediction within the historical
dataset; it does not directly assess forecasting for later
model years. The results should therefore not be assumed to
apply to modern vehicles or substantially different vehicle
populations.

Finally, the fitted coefficients describe associations within
the chosen models. Holding the other included predictors fixed
is a mathematical interpretation of a coefficient, not evidence
that changing a vehicle characteristic would cause the
corresponding change in fuel efficiency.

%We now want to discuss the mathods used to compare, evaluate and validate procedures. 
%Given any predictive, specify how data was split and which information was available during model selection
\subsection{Evaluation design}
\label{sec:evaluation}

Three approaches were evaluated: a baseline that predicts the
mean fuel efficiency of the training observations, a linear
regression using vehicle weight alone, and a multiple linear
regression using weight, horsepower and model year. The baseline
provides a reference for assessing whether vehicle characteristics
improve predictions beyond a single constant value.

\subsubsection{Training and test split}

The 392 complete observations were randomly divided into
313 training observations and 79 test observations, corresponding
approximately to an 80:20 split. A random seed of 42 was used
to make the split reproducible. All three approaches used the
same training and test observations, ensuring that differences
in performance were not due to different evaluation samples.

The split was random rather than chronological. Consequently,
the evaluation assesses prediction for held-out vehicles from
the historical dataset, rather than forecasting fuel efficiency
for later model years.

\subsubsection{Cross-validation and model comparison}

Five-fold cross-validation was performed within the training
set \cite[Section 5.1.3]{islr}. The 313 training observations
were shuffled and divided into five approximately equal groups,
using a random seed of 42. Each approach was fitted on four
groups and evaluated on the remaining group. This was repeated
five times so that each group served once as the validation set.

The same folds were used for all three approaches. For each
fold, the regression coefficients were estimated using only
that fold's training observations. Similarly, the baseline
prediction was calculated from the mean outcome in those
training observations, excluding the validation group.

Mean absolute error was calculated for each validation fold.
The mean and standard deviation of these five errors were
reported to describe average predictive performance and its
variation across folds. The standard deviation is descriptive;
it is not a confidence interval for predictive performance.

The scripts specify the predictor sets in advance of
cross-validation and do not perform an automated search over
alternative predictor combinations or tuning parameters.
Cross-validation therefore compares the three specified
approaches. It does not establish that the selected predictors
are the best possible subset of the available variables.

After cross-validation, each approach was fitted to all
313 training observations and evaluated on the same 79 test
observations. Test outcomes were used to calculate the final
evaluation metrics, not to estimate coefficients or calculate
cross-validation scores.

\subsubsection{Evaluation metrics}

Let $m$ denote the number of observations being evaluated,
$y_i$ the observed fuel efficiency, and $\widehat{y}_i$ its
prediction. Three complementary metrics were calculated.

\textbf{Mean absolute error.}
The mean absolute error measures the average magnitude of
prediction errors:
\begin{equation}
    \operatorname{MAE}
    = \frac{1}{m}\sum_{i=1}^{m}
      |y_i-\widehat{y}_i|.
    \label{eq:mae}
\end{equation}
It is expressed in miles per gallon, and smaller values
indicate better predictive performance.

\textbf{Root mean squared error.}
The root mean squared error is
\begin{equation}
    \operatorname{RMSE}
    = \sqrt{
        \frac{1}{m}\sum_{i=1}^{m}
        (y_i-\widehat{y}_i)^2
      }.
    \label{eq:rmse}
\end{equation}
It is also expressed in miles per gallon. Squaring the errors
before averaging makes this metric more sensitive to large
errors than mean absolute error. Smaller values indicate
better performance.

\textbf{Coefficient of determination.}
For an evaluation sample with nonzero outcome variation,
the coefficient of determination is
\begin{equation}
    R^2
    = 1-
      \frac{
        \sum_{i=1}^{m}(y_i-\widehat{y}_i)^2
      }{
        \sum_{i=1}^{m}(y_i-\overline{y}_{\mathrm{eval}})^2
      },
    \qquad
    \overline{y}_{\mathrm{eval}}
    = \frac{1}{m}\sum_{i=1}^{m}y_i.
    \label{eq:r-squared}
\end{equation}
This metric is dimensionless, and larger values indicate
better performance. A value of one represents perfect
predictions. A value of zero indicates the same squared error
as predicting the evaluation sample's mean, while a negative
value indicates greater squared error than that reference.

The evaluation mean in Equation~\ref{eq:r-squared} is used
only to calculate the metric. It differs from the training
mean used by the fitted baseline. Consequently, the baseline's
test-set $R^2$ is not necessarily zero.

Mean absolute error was used for cross-validation. All three
metrics were reported for the test set, alongside plots of
prediction errors to support interpretation.


\section{Results}
\label{sec:results}
\begin{table}[htbp]
\centering
\caption{Predictive performance on the same 79 test observations.
Mean absolute error and root mean squared error are measured
in miles per gallon; lower values indicate better performance.
The coefficient of determination ($R^2$) is dimensionless;
higher values indicate better performance.}
\label{tab:results}
\begin{tabularx}{\linewidth}{@{}X r r r@{}}
\toprule
\textbf{Method}
& \shortstack{\textbf{Mean absolute}\\\textbf{error}}
& \shortstack{\textbf{Root mean}\\\textbf{squared error}}
& $\boldsymbol{R^2}$ \\
\midrule
Mean baseline       & 5.881 & 7.185 & -0.011 \\
Weight regression   & 3.464 & 4.206 &  0.653 \\
Multiple regression & 2.509 & 3.255 &  0.792 \\
\bottomrule
\end{tabularx}
\end{table}

Table~\ref{tab:results} shows that the weight-only regression
reduced test mean absolute error from 5.881 to 3.464 miles
per gallon, an improvement of approximately 41.1\% relative
to the mean baseline. Its negative fitted weight coefficient
indicates that heavier vehicles were associated with lower
predicted fuel efficiency.

The multiple regression model reduced test mean absolute
error further, to 2.509 miles per gallon. This represents
an improvement of approximately 27.6\% compared with the
weight-only model. It also achieved the lowest mean
cross-validation error, at 2.695 miles per gallon, compared
with 3.317 for weight-only regression and 6.730 for the
mean baseline.

The agreement between the cross-validation and test-set
comparisons supports the conclusion that horsepower and
model year jointly add predictive value beyond vehicle
weight in this investigation. This comparison does not
establish the separate contribution of each additional
predictor or demonstrate causal effects on fuel efficiency.

\begin{figure}[htbp]
\centering
\includegraphics[width=\linewidth]{figures/regression_residual_comparison.png}
\caption{Test residuals against predicted fuel efficiency for
the weight-only and multiple regression models. Both panels
use the same axis scales. Residuals are observed minus predicted
fuel efficiency, measured in miles per gallon. Positive values
indicate underprediction, and negative values indicate
overprediction. The red dashed line marks zero prediction error.}
\label{fig:main-result}
\end{figure}


Figure~\ref{fig:main-result} illustrates the prediction errors
underlying the numerical comparison. The weight-only model
shows substantial variation in residuals at higher predicted
fuel efficiency. Although the multiple regression model
reduces overall prediction error, visible structure remains:
residuals tend to be positive at low predicted values and
negative across much of the middle range, with several
positive residuals at the upper end. This pattern suggests
that the additive linear specification may not fully capture
the relationship between the predictors and fuel efficiency.


%Here we interpret the findings and how they answeer the research question and/or compare to prvious work. Also practical implications
\newpage
\section{Discussion}
\label{sec:discussion}

Weight alone provided useful predictions, but including
horsepower and model year reduced errors further in both
cross-validation and the test set. This suggests that
vehicles with similar weights differ in ways that matter
for predicting fuel efficiency. Since horsepower and
model year were added together, this comparison does
not show how much each contributes separately.

Horsepower and model year may capture differences in engine
characteristics and vehicle design that weight alone misses.
This offers a possible explanation for the improvement,
although the analysis does not establish which mechanisms
are responsible.

The multiple regression model missed the observed test values by 2.5 miles per gallon on average in absolute terms, although some errors were much larger. No acceptable error threshold was defined for this project, so this result supports a comparison between models rather than a claim that the model is ready for practical use.

The lower average error does not mean that the multiple
regression model captured the relationship fully.
Patterns remain in the residuals, suggesting that some
prediction errors are systematic. These patterns give
a reason to explore other model specifications
\cite{fpp3-diagnostics}.


%Disscuss error patterns, sensitivity checks, unusual observations or alternative explanations relevant to this project.
\subsection{Diagnostics and robustness}

The residual plots show that prediction errors were not
uniformly distributed across the range of predicted fuel
efficiency. For the weight-only model, residuals were generally
positive at low predicted values, while a wider spread appeared
at higher predicted values. This indicates that the model
underpredicted fuel efficiency for several vehicles at the
lower end of its prediction range and produced substantial
errors in both directions at the upper end.

The multiple regression model reduced overall error, but
systematic patterns remained. Residuals were generally positive
at low predicted values and negative across much of the middle
range, with several positive residuals at the upper end.
This suggests that the additive linear specification may
not fully represent the relationship between the selected
predictors and fuel efficiency. Curvature, interactions or
omitted vehicle characteristics are possible explanations,
but the residual plots alone cannot distinguish between them
\cite{fpp3-diagnostics}.

Several test observations had comparatively large residuals.
These observations were retained: a large prediction error
does not itself establish that an observation is incorrect.
No formal influence analysis was performed, so the investigation
does not establish whether particular training observations
had a disproportionate effect on the fitted coefficients.

Five-fold cross-validation provided a limited assessment of
sensitivity to the training observations used for fitting.
The multiple regression model achieved a mean validation
absolute error of 2.695 miles per gallon, with a standard
deviation across folds of 0.243. The corresponding values
were 3.317 and 0.445 for weight-only regression, and 6.730
and 0.747 for the mean baseline.

The lower average error and smaller variation across folds
support the multiple regression model within this evaluation.
However, the fold standard deviations are descriptive and
do not constitute confidence intervals or a formal test
of differences between models. No repeated train--test splits,
alternative missing-value treatments or alternative predictor
specifications were evaluated. Robustness to those choices
therefore remains untested.



%Give a summary of main constraints and their consequences for interpretation. Consider data quality, sample size, assumptions, evaluation design and whether the findings generalise beyond the data studied
\subsection{Limitations}
The analysis used 392 complete observations, of which 79
were reserved for testing. Performance estimates therefore
depend on a relatively small test sample and one random
train--test split. Cross-validation provides additional
evidence from the training data, but does not remove uncertainty
about performance on other samples.

Six observations were excluded because horsepower values
were missing. Using the same complete observations allowed
consistent comparisons between models, but also removed
vehicles that could otherwise have been used in the weight-only
analysis. If the excluded vehicles differed systematically
from those retained, complete-case analysis could affect
the representativeness of the findings.

I explored the full cleaned dataset before fitting the
models. If those observations influenced my choice of
predictors, the test data indirectly informed model
development. The test set was excluded from fitting,
but was not reserved before all exploratory work.

Only two regression specifications and a constant baseline
were compared. The multiple regression model included weight,
horsepower and model year, but omitted other recorded
characteristics and contained no nonlinear or interaction
terms. Its superior performance among the evaluated approaches
does not establish that it is the best available model.
Furthermore, adding horsepower and model year together does
not identify their separate predictive contributions.

The remaining residual patterns indicate that some systematic
variation may be unexplained. No formal tests of error
variance, influence diagnostics or coefficient confidence
intervals were calculated. Consequently, the investigation
provides a predictive comparison rather than a comprehensive
assessment of statistical inference assumptions.

Finally, the dataset represents vehicles from model years
1970 to 1982. A random split evaluates prediction within
this historical sample, rather than forecasting for later
vehicle generations. The findings should not be assumed
to generalise to modern vehicles.


%Here we answer the research question, directly. Summaries all key findings and evidance that support this
\section{Conclusion}
\label{sec:conclusion}

Linear regression provided useful predictions of fuel
efficiency in the Auto MPG dataset. Weight-only regression
reduced test mean absolute error by 41.1\% compared with
the mean baseline. Adding horsepower and model year
reduced the error by a further 27.6\%, to 2.509 miles
per gallon, and also gave the lowest mean
cross-validation error.

Nevertheless, systematic residual patterns remained, and
the evaluation was limited to a small historical dataset.
The multiple regression model was the strongest of the three
approaches tested, but its performance should not be assumed
to extend to modern vehicles or interpreted as evidence
of causal relationships.


%Add a short section to identify a few possible next steps and why they would aid investigation
\subsection{Future work}
A first extension would investigate nonlinear and interaction
terms, guided by the residual patterns. For example, a squared
weight term could represent curvature, while an interaction
between weight and horsepower could allow their predictive
relationship to vary jointly. These extensions should be
compared using validation data to determine whether the
additional flexibility improves prediction.

A second step would assess sensitivity to modelling choices.
Comparing models that add horsepower and model year separately
would help distinguish their predictive contributions.
Alternative treatments of missing horsepower values could
also show whether excluding incomplete observations materially
affects the findings. Any estimated replacement values should
be learned within each training fold to avoid using validation
information.

Finally, further model development should use a clearly
separated evaluation procedure. Repeated or nested
cross-validation could assess sensitivity to data partitions
and account for model selection. A new, untouched dataset
would provide stronger evidence of generalisation; evaluating
on more recent vehicles would specifically test whether
the relationships extend beyond the historical sample.

% --- References -------------------------------------------------------
% Simple built-in bibliography: no separate .bib file is required.
% Replace this example entry with real source details before publishing.
% Add citations in the report using \cite{source-key}.
\clearpage
\phantomsection
\addcontentsline{toc}{section}{References}
\begin{thebibliography}{9}
\bibitem{auto-mpg} %Auto-Mpg data
Quinlan, R. (1993).
\textit{Auto MPG} [Dataset].
UCI Machine Learning Repository.
\url{https://doi.org/10.24432/C5859H}.

\bibitem{illinois-least-squares}
CS 357 Course Staff.
\textit{Least Squares Fitting}.
University of Illinois Urbana-Champaign, CS 357 Textbook.
See the section ``Normal Equations''.
\url{https://cs357.cs.illinois.edu/textbook/notes/linear-least-squares.html}.
Accessed 29 September 2026.

\bibitem{iowa-mean-proof}
\textit{Basic Statistics}.
Course notes hosted by Iowa State University.
Section 2.1, Theorem 1(a).
\url{https://www2.econ.iastate.edu/classes/econ500/hallam/documents/Basic_stat.pdf}.

\bibitem{penn-state-nested-models}
Pennsylvania State University.
\textit{The General Linear F-Test}.
STAT 462, Section 5.6.
\url{https://online.stat.psu.edu/stat462/node/135/}.
Accessed 29 September 2026.

\bibitem{islr}
James, G., Witten, D., Hastie, T. and Tibshirani, R. (2013).
\textit{An Introduction to Statistical Learning:
with Applications in R}.
Springer.
See Section 5.1.3 for $k$-fold cross-validation.
\url{https://www.statlearning.com/s/ISLRSeventhPrinting.pdf}.

\bibitem{fpp3-diagnostics}
Hyndman, R. J. and Athanasopoulos, G. (2021).
\textit{Forecasting: Principles and Practice}.
3rd edition. OTexts, Melbourne, Australia.
Section 7.3, ``Evaluating the regression model''.
\url{https://otexts.com/fpp3/regression-evaluation.html}.
Accessed 29 September 2026.
\end{thebibliography}
% --- Optional appendix -----------------------------------------------
\clearpage
\appendix

\end{document}
